\documentclass[11pt]{amsart}
\usepackage[margin=1in]{geometry}
\usepackage{amsmath,amssymb,amsthm,mathtools}
\usepackage[T1]{fontenc}
\usepackage{lmodern}
\usepackage{microtype}
\usepackage{enumitem}
\usepackage[hidelinks]{hyperref}

\newtheorem{theorem}{Theorem}[section]
\newtheorem{lemma}[theorem]{Lemma}
\newtheorem{proposition}[theorem]{Proposition}
\newtheorem{corollary}[theorem]{Corollary}
\theoremstyle{definition}
\newtheorem{definition}[theorem]{Definition}
\newtheorem{example}[theorem]{Example}
\theoremstyle{remark}
\newtheorem{remark}[theorem]{Remark}

\newcommand{\R}{\mathbb{R}}
\DeclareMathOperator{\JM}{J\!M}

\title[Topological Frustration]{Topological Frustration and Residual Cost:\\
Exact Local Minimizers for Reciprocal Bond Geometry}

\author{Jonathan Washburn}
\address{Austin, Texas, USA}
\email{washburn.jonathan@gmail.com}
\date{February 2026}

\begin{document}
\begin{abstract}
Let $J(x)=\frac12(x+x^{-1})-1$ for $x>0$.
For fixed positive neighbor data $n_1,\dots,n_d$, we study the local bond
cost $\mathcal C(v)=\sum_{i=1}^d J(v/n_i)$.
We prove an exact closed form for the unique minimizer:
\[
v^\star=\sqrt{\frac{\sum_i n_i}{\sum_i n_i^{-1}}}
=:\JM(n_1,\dots,n_d),
\]
and an exact minimal residual
$\min_{v>0}\mathcal C(v)=\sqrt{(\sum_i n_i)(\sum_i n_i^{-1})}-d$.
This residual is strictly positive unless all neighbors are equal
(by Cauchy--Schwarz).
For two incompatible neighbors $n_1\neq n_2$, we obtain the explicit
frustration formula
$\min_{v>0}(J(v/n_1)+J(v/n_2))=\sqrt{n_1/n_2}+\sqrt{n_2/n_1}-2>0$.
At graph level, if local neighborhoods produce different $J$-means at two
vertices, then any simultaneous local minimizer field is necessarily
non-uniform.  We call this forced non-uniformity \emph{topological
frustration} and prove it in the precise form:
the optimizing field at each node is uniquely determined by its
neighborhood (strict convexity), so distinct neighborhoods force distinct
node values.  We record application-interface statements for fluid-direction
smoothing, constraint relaxation, and lattice embedding.
\end{abstract}

\subjclass[2020]{Primary 26A51; Secondary 05C69, 49J40, 26D15}
\keywords{Reciprocal cost, strict frustration, local minimizer,
harmonic-arithmetic mean, non-collapse, graph optimization}
\maketitle

%=======================================================================
\section{Introduction}\label{sec:intro}
%=======================================================================

The reciprocal cost $J(x)=\frac12(x+x^{-1})-1$ is nonnegative,
reciprocal-symmetric ($J(x)=J(1/x)$), and strictly convex on $(0,\infty)$.
These one-variable facts are developed in detail in the companion
paper~\cite{F1}, where it is shown that $J$ is the unique such function
satisfying the d'Alembert functional identity with unit curvature.

This paper addresses the \emph{multi-bond local problem}:
given fixed bond neighbor values $n_1,\dots,n_d>0$,
what value $v>0$ minimizes $\sum_{i=1}^d J(v/n_i)$?
The answer is an explicit closed form.  This immediately yields precise
frustration inequalities: when neighbors disagree, the residual cost is
provably positive, with an exact formula.

The paper is organized as follows.
Section~\ref{sec:local} derives the closed form and the $J$-mean.
Section~\ref{sec:frustration} proves the strict frustration inequalities
with explicit lower bounds.
Section~\ref{sec:graph} lifts the local result to a graph-level
non-collapse theorem.
Section~\ref{sec:apps} records application interfaces used by companion
papers.

%=======================================================================
\section{Local optimization against fixed neighbors}\label{sec:local}
%=======================================================================

\begin{definition}[Local bond cost]
For fixed $n_1,\dots,n_d>0$ ($d\ge1$), define
\[
\mathcal C(v):=\sum_{i=1}^d J(v/n_i),\qquad v>0.
\]
\end{definition}

\begin{lemma}[Closed form]\label{lem:closed}
With $S_1:=\sum_{i=1}^d n_i$ and $S_{-1}:=\sum_{i=1}^d n_i^{-1}$,
\begin{equation}\label{eq:closed}
\mathcal C(v)=\frac12\!\left(S_{-1}\,v+\frac{S_1}{v}\right)-d.
\end{equation}
\end{lemma}
\begin{proof}
Each summand expands as
$J(v/n_i)=\frac12(v/n_i+n_i/v)-1$.
Summing over $i$:
\[
\mathcal C(v)
=\frac{v}{2}\sum_i\frac{1}{n_i}+\frac{1}{2v}\sum_i n_i-d
=\frac12\!\left(S_{-1}\,v+\frac{S_1}{v}\right)-d.
\qedhere
\]
\end{proof}

\begin{theorem}[Exact minimizer and strict convexity]\label{thm:min}
$\mathcal C$ is strictly convex on $(0,\infty)$ with unique minimizer
\begin{equation}\label{eq:vstar}
v^\star=\sqrt{\frac{S_1}{S_{-1}}}
=\sqrt{\frac{\sum_i n_i}{\sum_i n_i^{-1}}}.
\end{equation}
\end{theorem}
\begin{proof}
Differentiate~\eqref{eq:closed}:
\[
\mathcal C'(v)=\frac12\!\left(S_{-1}-\frac{S_1}{v^2}\right),\qquad
\mathcal C''(v)=\frac{S_1}{v^3}>0.
\]
Strict positivity of $\mathcal C''$ gives strict convexity.
Since $\mathcal C(v)\to+\infty$ as $v\to0^+$ or $v\to+\infty$
(from~\eqref{eq:closed}), $\mathcal C$ is coercive.
A strictly convex coercive function has exactly one critical point,
which is the global minimum.
Setting $\mathcal C'(v)=0$: $S_{-1}=S_1/v^2\implies v^2=S_1/S_{-1}
\implies v^\star=\sqrt{S_1/S_{-1}}$.
\end{proof}

\begin{definition}[$J$-mean]\label{def:jmean}
For positive data $n_1,\dots,n_d$, define the \emph{$J$-mean} as
\[
\JM(n_1,\dots,n_d):=\sqrt{\frac{\sum_i n_i}{\sum_i n_i^{-1}}}.
\]
\end{definition}

\begin{corollary}[AM/HM interpretation]\label{cor:means}
With $\operatorname{AM}:=\frac1d\sum_i n_i$ and
$\operatorname{HM}:=d/\sum_i n_i^{-1}$,
\[
\JM(n_1,\dots,n_d)=\sqrt{\operatorname{AM}\cdot\operatorname{HM}}.
\]
For $d=2$: $\JM(n_1,n_2)=\sqrt{n_1n_2}$ (the geometric mean).
\end{corollary}
\begin{proof}
$S_1/S_{-1}=d\cdot\operatorname{AM}\cdot\operatorname{HM}/d
=\operatorname{AM}\cdot\operatorname{HM}$.
For $d=2$:
$(n_1+n_2)/(n_1^{-1}+n_2^{-1})
=(n_1+n_2)\cdot n_1n_2/(n_1+n_2)=n_1n_2$.
\end{proof}

\begin{remark}[Relation to other means]
The $J$-mean satisfies $\operatorname{HM}\le\JM\le\operatorname{AM}$
by the AM-HM inequality, with equality throughout iff all $n_i$ are equal.
It coincides with the geometric mean for $d=2$ but differs for $d\ge3$
(for instance, $\JM(1,1,4)=\sqrt{6/(1+1+1/4)}=\sqrt{6/(9/4)}
=\sqrt{8/3}\approx1.633$, while $\operatorname{GM}(1,1,4)=4^{1/3}\approx1.587$).
See Bullen~\cite{Bullen} for a comprehensive discussion of mean inequalities.
\end{remark}

\begin{proposition}[Exact minimum value]\label{prop:min-value}
\[
\min_{v>0}\mathcal C(v)=\sqrt{S_1 S_{-1}}-d\ge0.
\]
Equality holds iff $n_1=\cdots=n_d$.
\end{proposition}
\begin{proof}
Evaluate~\eqref{eq:closed} at $v^\star=\sqrt{S_1/S_{-1}}$:
\begin{align*}
\mathcal C(v^\star)
&=\frac12\!\left(S_{-1}\sqrt{S_1/S_{-1}}
  +\frac{S_1}{\sqrt{S_1/S_{-1}}}\right)-d \\
&=\frac12(\sqrt{S_1 S_{-1}}+\sqrt{S_1 S_{-1}})-d
=\sqrt{S_1 S_{-1}}-d.
\end{align*}
By the Cauchy--Schwarz inequality in $\R^d$
(applied to the vectors $(n_1^{1/2},\dots,n_d^{1/2})$ and
$(n_1^{-1/2},\dots,n_d^{-1/2})$):
\[
S_1 S_{-1}
=\Bigl(\sum_i n_i\Bigr)\Bigl(\sum_i n_i^{-1}\Bigr)
\ge\Bigl(\sum_i 1\Bigr)^2=d^2.
\]
Hence $\sqrt{S_1 S_{-1}}\ge d$, giving $\mathcal C(v^\star)\ge0$.
Equality in Cauchy--Schwarz occurs iff $n_i^{1/2}$ and $n_i^{-1/2}$
are proportional for all $i$, i.e.\ iff all $n_i$ are equal.
\end{proof}

\begin{example}
For $d=3$, $n_1=1$, $n_2=2$, $n_3=8$:
$S_1=11$, $S_{-1}=1+1/2+1/8=13/8$,
$v^\star=\sqrt{88/13}\approx 2.602$,
$\min\mathcal C=\sqrt{143/8}-3\approx 1.227$.
\end{example}

%=======================================================================
\section{Strict frustration inequalities}\label{sec:frustration}
%=======================================================================

The central message of this section: when neighbor data disagree,
no choice of $v$ can bring the total bond cost to zero.

\begin{theorem}[No simultaneous annihilation]\label{thm:no-simult}
If $n_1,n_2>0$ and $n_1\neq n_2$, there is no $v>0$ with
$J(v/n_1)=J(v/n_2)=0$.
\end{theorem}
\begin{proof}
By~\cite[Corollary~2.5]{F1}, $J(v/n_i)=0\iff v=n_i$.
Simultaneous vanishing forces $v=n_1$ and $v=n_2$, contradicting
$n_1\neq n_2$.
\end{proof}

\begin{corollary}[Strict two-bond frustration]\label{cor:strict-two}
If $n_1\neq n_2$ ($n_1,n_2>0$), then for every $v>0$:
$J(v/n_1)+J(v/n_2)>0$.
\end{corollary}
\begin{proof}
Each summand is $\ge0$ (since $J\ge0$).
If the sum were zero, both summands would vanish,
contradicting Theorem~\ref{thm:no-simult}.
\end{proof}

\begin{proposition}[Exact two-neighbor residual]\label{prop:two-exact}
For $n_1,n_2>0$:
\begin{equation}\label{eq:two-residual}
\min_{v>0}\bigl(J(v/n_1)+J(v/n_2)\bigr)
=\sqrt{\frac{n_1}{n_2}}+\sqrt{\frac{n_2}{n_1}}-2.
\end{equation}
If $n_1\neq n_2$, this minimum is strictly positive.
\end{proposition}
\begin{proof}
By Corollary~\ref{cor:means}, the minimizer is $v^\star=\sqrt{n_1n_2}$.
Set $r:=\sqrt{n_2/n_1}$, so $v^\star/n_1=r$ and $v^\star/n_2=1/r$.
Then:
\begin{align*}
J(r)+J(1/r)
&=\frac12(r+r^{-1})-1+\frac12(r^{-1}+r)-1 \\
&=(r+r^{-1})-2 \\
&=\sqrt{\frac{n_2}{n_1}}+\sqrt{\frac{n_1}{n_2}}-2.
\end{align*}
By the AM-GM inequality, $r+r^{-1}\ge 2\sqrt{r\cdot r^{-1}}=2$,
with equality iff $r=1$ (i.e.\ $n_1=n_2$).
\end{proof}

\begin{corollary}[$d$-bond positivity from one incompatible pair]
\label{cor:d-bond}
Let $n_1,\dots,n_d>0$ with $n_i\neq n_j$ for some pair $i\neq j$.
Then for every $v>0$:
\[
\sum_{\ell=1}^d J(v/n_\ell)>0.
\]
Moreover,
\[
\sum_{\ell=1}^d J(v/n_\ell)
\ge\sqrt{\frac{n_i}{n_j}}+\sqrt{\frac{n_j}{n_i}}-2>0.
\]
\end{corollary}
\begin{proof}
All $d$ terms are $\ge0$.  Drop all except the $(i,j)$ pair:
$\sum_\ell J(v/n_\ell)\ge J(v/n_i)+J(v/n_j)$.
Apply Corollary~\ref{cor:strict-two} for strict positivity and
Proposition~\ref{prop:two-exact} for the explicit bound
(noting that the bound holds for every $v$, hence also at the $d$-bond
optimum).
\end{proof}

\begin{example}
With $n_1=2$, $n_2=8$: residual $=\sqrt{1/4}+\sqrt{4}-2=1/2+2-2=1/2$.
With $n_1=1$, $n_2=100$: residual $=1/10+10-2=8.1$.
Wide disparity produces large frustration.
\end{example}

%=======================================================================
\section{Graph-level non-collapse}\label{sec:graph}
%=======================================================================

We now lift the local frustration to a global statement on graphs.

\begin{definition}[Prescribed neighborhood data]\label{def:prescribed}
Let $G=(V,E)$ be a finite connected graph.
For each vertex $u\in V$, fix a positive multiset
$\mathcal N(u)=\{n_{u,1},\dots,n_{u,d_u}\}$.
The \emph{local $J$-mean} is $m_u:=\JM(\mathcal N(u))$.
\end{definition}

\begin{theorem}[Distinct local means force distinct optimal states]
\label{thm:distinct}
Suppose there exist vertices $u,w\in V$ with $m_u\neq m_w$.
If a field $x=(x_v)_{v\in V}$ satisfies
\[
x_u\in\arg\min_{t>0}\sum_{i=1}^{d_u}J(t/n_{u,i})
\quad\text{for every }u\in V,
\]
then $x_u\neq x_w$.  In particular, no such field is spatially uniform.
\end{theorem}
\begin{proof}
By Theorem~\ref{thm:min}, each local $\arg\min$ is a singleton:
$\arg\min_{t>0}\sum_i J(t/n_{u,i})=\{m_u\}$.
Hence $x_u=m_u$ for every $u$.  If $m_u\neq m_w$, then $x_u\neq x_w$.
A uniform field $x_u\equiv c$ would require $m_u=c$ for all $u$,
contradicting $m_u\neq m_w$.
\end{proof}

\begin{remark}[Scope of the non-collapse statement]\label{rem:scope}
Theorem~\ref{thm:distinct} concerns optimization against
\emph{prescribed} (externally fixed) neighborhood data.
For the unconstrained edge-energy functional
$\sum_{\{u,w\}\in E}J(x_u/x_w)$, constant fields achieve zero energy
and are trivially optimal.
Therefore the topological frustration mechanism requires either
incompatible prescribed boundary conditions, external forcing, or
additional constraint structure (such as window neutrality~\cite{F3}).
This distinction is important in applications.
\end{remark}

%=======================================================================
\section{Application interfaces}\label{sec:apps}
%=======================================================================

We record the precise forms used by companion papers.

\begin{remark}[Fluid-direction smoothing interface]\label{rem:fluid}
If direction updates near the vorticity zero set $\{\omega=0\}$ are
modeled by local reciprocal cost minimization against neighboring
directions, then incompatible local neighborhoods generate strictly
positive residual cost (Corollary~\ref{cor:d-bond}), preventing
collapse to a single linearized average in defect regions.
See~\cite{F6} for the full blow-up exclusion argument.
\end{remark}

\begin{remark}[Constraint-relaxation interface]\label{rem:constraint}
In graph encodings of combinatorial constraints, if a relaxed local
state must satisfy incompatible neighboring targets, the strict residual
cost of Corollary~\ref{cor:strict-two} forbids a zero-cost fractional
survivor.  See the companion paper on combinatorial feasibility for the
complete reduction.
\end{remark}

\begin{remark}[Lattice embedding interface]\label{rem:lattice}
In reciprocal-cost lattice updates, distinct local prescribed
neighborhoods force differentiated local optima by
Theorem~\ref{thm:distinct}, providing a direct non-collapse criterion
for gauge-field lattice embeddings.
\end{remark}

%=======================================================================
\section*{Acknowledgments}
%=======================================================================
The author thanks the referees for careful reading.

\begin{thebibliography}{99}

\bibitem{Bullen}
P.~S.~Bullen,
\emph{Handbook of Means and Their Inequalities},
Mathematics and Its Applications, vol.~560, Kluwer, 2003.

\bibitem{F1}
J.~Washburn,
\emph{The canonical reciprocal cost: exact multi-bond minimization,
frustration lower bounds, and simultaneous-vs-sequential descent},
preprint, 2026.

\bibitem{F3}
J.~Washburn,
\emph{Window neutrality, future-balance projection, and the phantom
balance constraint},
preprint, 2026.

\bibitem{F6}
J.~Washburn,
\emph{Alexander-duality linking, link penalties, and the finite-capacity
veto in three dimensions},
preprint, 2026.

\end{thebibliography}

\end{document}
